\documentclass[11pt,a4paper]{article}
\usepackage[utf8]{inputenc}
\usepackage[portuguese]{babel}
\usepackage[T1]{fontenc}
\usepackage{amsmath,amssymb}
\usepackage{graphicx}
\usepackage{booktabs}
\usepackage{hyperref}
\usepackage{geometry}
\usepackage{float}
\usepackage{listings}
\usepackage{enumerate}

\geometry{verbose,tmargin=2.5cm,bmargin=2.5cm,lmargin=3cm,rmargin=3cm}

\title{Análise de Valor no Mercado Automóvel\\
\large Desenvolvimento de um Sistema Integrado Local e Web com Machine Learning para Avaliação de Carros Usados}

\author{Juliana Rodrigues (1231182) \and Geovani Junior (1131518)\\
\small Instituto Superior de Engenharia do Porto}
\date{Versão Final --- Junho de 2026}

\begin{document}

\maketitle

\begin{abstract}
O presente trabalho descreve o desenvolvimento de um sistema integrado para a recolha automática, a análise de dados e a comercialização de anúncios de veículos usados provenientes do mercado online OLX Portugal. A solução assenta numa arquitetura de componentes integrados: (1) um scraper em Python (\texttt{olx\_car\_scraper.py}) que recolhe dados estruturados de anúncios de automóveis, incluindo marca, modelo, ano, quilometragem, preço, potência e localização; (2) um módulo de inteligência artificial (\texttt{gerar\_previsoes\_ia.py}) que utiliza algoritmos de aprendizagem automática, nomeadamente XGBoost, Random Forest, LightGBM e modelos de regressão linear/exponencial, para estimar o valor justo de mercado de cada veículo com base nas suas características; (3) um workflow de automação construído na plataforma n8n, responsável pela filtragem dos anúncios já avaliados e pelo envio de alertas via Telegram sempre que surge um veículo abaixo do valor estimado; (4) um portal web de comercialização desenvolvido em HTML5, CSS3 e JavaScript vanilla, com filtros avançados e indicadores visuais de oportunidade; e (5) uma aplicação desktop desenvolvida em C++ com o framework Qt (\texttt{GeOLXavani}) que permite a gestão local do pipeline, treino de modelos e avaliação personalizada através de scripts auxiliares. O sistema evidencia o potencial da unificação de dados e modelos preditivos persistidos em formato padrão para mitigar a assimetria de informação no mercado automóvel nacional.

\textbf{Keywords:} Web Scraping, Machine Learning, XGBoost, Random Forest, LightGBM, n8n, Qt C++, Mercado Automóvel, OLX, Portal Web, Python, Data Science.
\end{abstract}

\section{Introdução}
O mercado de automóveis usados em Portugal movimenta anualmente centenas de milhares de transações. Plataformas como o OLX Portugal concentram uma parte significativa deste volume, com milhares de novos anúncios publicados diariamente.
A elevada volatilidade dos preços, a depreciação não linear dos veículos e o elevado número de listagens tornam difícil, para um comprador individual, identificar manualmente as oportunidades com melhor relação qualidade/preço.

Visando mitigar a assimetria de informação existente, este projeto propõe o desenvolvimento de uma arquitetura unificada, baseada em automação, Inteligência Artificial e interfaces local e web. A primeira camada é responsável pela extração sistemática de dados brutos do mercado. A segunda foca-se no pré-processamento de dados e no treino de modelos preditivos para estimar o valor comercial justo de cada automóvel. Por sua vez, a terceira camada orquestra a análise destas avaliações, emitindo alertas em tempo real perante oportunidades de negócio vantajosas. Por fim, o ecossistema disponibiliza interfaces duplas: um portal Web interativo, que centraliza esta inteligência analítica para apoiar a decisão final do consumidor, e uma aplicação desktop nativa desenvolvida em C++/Qt que permite o controlo do scraper, retreino de modelos e uma ferramenta de avaliação manual de veículos offline baseada em modelos persistidos.

\section{Arquitetura do Sistema}
O sistema é composto por módulos interdependentes que formam um pipeline de dados completo. O fluxo de execução é orquestrado de forma a garantir a passagem sequencial e íntegra dos dados entre a extração, o processamento preditivo, a interface visual e as ferramentas desktop.

\subsection{Módulo de Recolha de Dados (\texttt{olx\_car\_scraper.py})}
O primeiro módulo do sistema é responsável pela aquisição automatizada de dados reais e atualizados sobre o mercado automóvel. Para o efeito, desenvolveu-se um scraper em Python direcionado à plataforma OLX. O componente opera em modo headless por omissão, mimetiza o comportamento de um utilizador humano através de pausas aleatórias entre pedidos (1,8--4,2 segundos) e rotação de user-agents, e aceita automaticamente os banners de consentimento de cookies. Durante o parsing dos dados detalhados, o scraper extrai também atributos específicos de hardware, tais como a potência em cavalos (\texttt{potencia\_cv}) a partir dos campos do anúncio.

\subsection{Filtros de Ano e Preço}
A função \texttt{passa\_filtros()} implementa lógica de filtragem em dois momentos: após o parsing do cartão de listagem (filtragem rápida, evita visitas desnecessárias à página de detalhe) e após a visita ao detalhe (filtragem final com dados completos). A função aceita limites mínimos e máximos para ano e preço, tratando valores \texttt{None} como ``sem restrição'' para não descartar anúncios onde o campo ainda não foi extraído.

\section{Módulo de Previsão de Preços por IA (\texttt{gerar\_previsoes\_ia.py})}
Este módulo atua como o motor analítico do projeto. O seu objetivo principal é processar os dados brutos recolhidos no portal OLX, treinar múltiplos algoritmos de Machine Learning em simultâneo e prever o valor justo de mercado para cada veículo. O resultado final enriquece a base de dados do front-end, fornecendo ao utilizador uma avaliação imediata sobre a qualidade do negócio. Para garantir a precisão e robustez do sistema, o desenvolvimento deste módulo foi dividido em três fases fundamentais: preparação de dados, modelagem preditiva e persistência física.

\subsection{Preparação e Limpeza de Dados}
Antes do treino da Inteligência Artificial, os dados extraídos pelo scraper passam por um rigoroso processo de limpeza. Esta fase inclui:
\begin{itemize}
    \item \textbf{Tratamento de Anomalias:} Remoção de anúncios sem preço definido ou com valores fora dos limites plausíveis de mercado (ex. quilometragens irrealistas ou preços extremos).
    \item \textbf{Imputação de Dados Faltantes:} Quando atributos cruciais, como os quilómetros ou o ano de fabrico, não são detetados, o sistema utiliza a mediana global do conjunto de dados para preencher a lacuna.
    \item \textbf{Estimação de Potência por Heurísticas:} Caso a potência (\texttt{hp}) não tenha sido capturada no scraper, o sistema aplica um processamento de texto baseado em termos frequentes no título (como ``320d'', ``1.6 tdi'') e a marca do veículo para atribuir uma potência estimada razoável, eliminando a anterior limitação de atribuir um valor flat de 100~cv.
    \item \textbf{Engenharia de Features:} Conversão de variáveis categóricas (Marca, Modelo simplificado, Combustível e Transmissão) em codificação dummy (\texttt{get\_dummies}) e criação de novas variáveis temporais como a idade do carro (\texttt{age}) e a média de quilómetros percorridos por ano (\texttt{km\_per\_year}). Para os modelos lineares, as features numéricas passam por uma normalização estatística (\texttt{StandardScaler}).
\end{itemize}

\subsection{Modelos Utilizados}
Foram avaliados seis modelos de regressão, cobrindo desde abordagens lineares simples até métodos de ensemble baseados em árvores de decisão. A descrição metodológica de cada modelo está detalhada na Tabela~\ref{tab:modelos_desc}.

\begin{table}[H]
\centering
\caption{Modelos de regressão avaliados.}
\label{tab:modelos_desc}
\begin{tabular}{ll}
\toprule
\textbf{Modelo} & \textbf{Descrição} \\
\midrule
Regressão Linear & Modelo de referência (baseline) com relação linear direta entre variáveis. \\
Regressão Exponencial & Aplica o logaritmo ao target ($y$), captando a desvalorização percentual. \\
Regressão Logarítmica & Aplica logaritmo às features numéricas, modelando o efeito de quilómetros. \\
Random Forest & Conjunto de árvores de decisão que captura não-linearidades do mercado. \\
XGBoost & Gradient boosting com regularização, focado na otimização de interações. \\
LightGBM & Gradient boosting otimizado para rapidez e árvores orientadas a folhas. \\
\bottomrule
\end{tabular}
\end{table}

\subsection{Métricas de Avaliação}
A escolha do modelo final assenta principalmente no MAPE (Mean Absolute Percentage Error), em detrimento do MAE absoluto. Esta opção justifica-se pela natureza do dataset, que cobre veículos com preços muito díspares: um erro absoluto de, por exemplo, 1000 EUR  representa um desvio relativo muito diferente consoante o veículo custe 5.000 EUR  ou 50.000 EUR. O MAPE normaliza esta diferença, permitindo uma comparação mais interpretável entre modelos. As restantes métricas (MAE, RMSE, $R^2$) são reportadas como complemento.

A Tabela~\ref{tab:metricas_exec} apresenta as métricas resultantes obtidas na execução final sobre o conjunto de teste unificado de 501 anúncios reais.

\begin{table}[H]
\centering
\caption{Métricas de avaliação obtidas por modelo.}
\label{tab:metricas_exec}
\begin{tabular}{lrrrr}
\toprule
\textbf{Modelo} & \textbf{MAE (EUR)} & \textbf{RMSE (EUR)} & \textbf{$R^2$} & \textbf{MAPE (\%)} \\
\midrule
Regressão Linear & 10.237,2 & 17.813,4 & 0,612 & 108,9\% \\
Regressão Exponencial & \textbf{8.411,0} & \textbf{17.500,1} & 0,625 & \textbf{51,8\%} \\
Regressão Logarítmica & 9.690,2 & 17.422,6 & \textbf{0,629} & 84,2\% \\
Random Forest & 11.172,2 & 21.908,3 & 0,413 & 115,5\% \\
XGBoost & 11.016,5 & 20.604,8 & 0,481 & 93,7\% \\
LightGBM & 10.963,3 & 21.003,5 & 0,460 & 114,5\% \\
\bottomrule
\end{tabular}
\end{table}

O modelo baseado em \textbf{Regressão Exponencial (Ridge)} apresentou o melhor compromisso global, registando o menor erro percentual médio (MAPE de 51,8\%) e o menor desvio absoluto médio (MAE de 8.411,0 EUR), sendo selecionado para a produção e persistido.

\subsection{Gráficos Comparativos}
O script exporta a figura de comparação de modelos (\texttt{comparacao\_modelos.png}) contendo gráficos de dispersão (preço previsto vs real) e comparadores de barras para as métricas, permitindo diagnosticar o overfitting de algoritmos com elevada flexibilidade como o Random Forest no nosso dataset atual.

\subsection{Do Terminal para Produção}
A transição da fase exploratória para a fase de produção traduz-se no uso do script \texttt{gerar\_previsoes\_ia.py}, executado em modo batch sem interação do utilizador. Este processa o dataset unificado completo e grava as previsões resultantes no ficheiro \texttt{olx\_carros\_ia.json}, integrado no portal web através dos campos \texttt{preco\_ia} e \texttt{status\_ia}.

\subsection{Integração com o Pipeline}
Ao contrário de uma integração direta no workflow n8n, a chamada ao módulo de IA é feita pelo script unificador \texttt{executar\_tudo.py}, que invoca sequencialmente o scraper e o módulo de previsão. O n8n atua após a conclusão deste pipeline, lendo o ficheiro já enriquecido e aplicando a lógica de deteção de oportunidades.

\section{Workflow de Automação com n8n}
O n8n opera como camada de orquestração assíncrona dos alertas: é executado de forma agendada, lê o ficheiro final e notifica o utilizador caso sejam identificados desvios favoráveis. O fluxo proposto é composto por nove nós funcionais, descritos na Tabela~\ref{tab:n8n_nodes}.

\begin{table}[H]
\centering
\caption{Nós do workflow n8n e respetivas funções.}
\label{tab:n8n_nodes}
\begin{tabular}{lll}
\toprule
\textbf{\#} & \textbf{Nó} & \textbf{Função} \\
\midrule
1 & Schedule Trigger & Execução periódica agendada (e.g. a cada 6 horas). \\
2 & Execute Command & Executa o script \texttt{executar\_tudo.py} para iniciar o pipeline. \\
3 & Read/Write Files & Lê o ficheiro final gerado pelo previsor de IA. \\
4 & Spreadsheet File & Converte o ficheiro em registos estruturados. \\
5 & Filter / IF & Aplica filtros básicos de preço, ano e marca. \\
6 & DB Lookup & Verifica se o anúncio já foi processado e armazenado. \\
7 & DB Insert & Guarda novos anúncios na base de dados (Airtable/Supabase). \\
8 & Code (JS) & Formata a mensagem de alerta com os dados detalhados do veículo. \\
9 & Switch + Telegram & Envia a notificação com os desvios calculados. \\
\bottomrule
\end{tabular}
\end{table}

\subsection{Deteção de Oportunidades}
O nó de formatação calcula o desvio de cada anúncio relativamente ao preço de referência, que corresponde à previsão da IA (\texttt{preco\_ia}) quando disponível, ou à mediana histórica como fallback. Anúncios com preço inferior à referência em mais de 10\% são classificados como ``oportunidade'' e desencadeiam o envio do alerta via Telegram.

\section{Portal Web de Comercialização}
O portal AutoMercado é uma aplicação web de página única (SPA) desenvolvida em HTML5, CSS3 e JavaScript vanilla. A interface oferece uma barra de pesquisa rápida, filtros avançados (preço, ano, quilometragem, combustível e transmissão), ordenação dinâmica e um sistema de favoritos com persistência local.

\subsection{Indicadores de Preço}
Cada card de veículo exibe o desvio em relação ao valor justo estimado pela IA. O desvio é realçado visualmente a verde (com a poupança associada) se o valor pedido estiver abaixo do previsto pela IA, ou a vermelho caso contrário, orientando intuitivamente as decisões de compra.

\section{Aplicação Desktop Qt (\texttt{GeOLXavani})}
Para além do portal web, o ecossistema do projeto foi expandido para incluir uma aplicação desktop desenvolvida em C++ com o framework Qt (Qt Widgets). Esta aplicação (\texttt{GeOLXavani}) destina-se a utilizadores que necessitam de uma gestão local mais robusta, oferecendo controlos diretos sobre a recolha de dados, treino de modelos e avaliação de veículos individualmente.

A interface está estruturada em quatro widgets principais interligados através de um controlador central (\texttt{MainWindow}):
\begin{itemize}
    \item \textbf{\texttt{ScraperWidget}:} Permite controlar a execução do script Python de web scraping através de uma interface visual. O utilizador pode configurar parâmetros como o número de páginas a extrair e ver em tempo real o log de execução do terminal diretamente numa caixa de texto.
    \item \textbf{\texttt{MlWidget}:} Oferece controlos para iniciar o treino local dos modelos de regressão a partir da interface gráfica, lendo o ficheiro de dados unificado e reportando as métricas de performance.
    \item \textbf{\texttt{TrainingWidget}:} Exibe as métricas comparativas resultantes do treino dos modelos e permite despoletar a geração de previsões em lote (\texttt{gerar\_previsoes\_ia.py}), atualizando o ficheiro JSON e gravando o modelo preditivo final.
    \item \textbf{\texttt{CarSearchWidget}:} Apresenta uma tabela interativa com todos os veículos registados no ficheiro unificado. Possui opções de filtragem e pesquisa por marca ou modelo, e integra um botão para avaliação personalizada do veículo.
\end{itemize}

A avaliação personalizada é executada no widget \texttt{CarEvaluationWidget}, o qual invoca o script independente \texttt{evaluate\_car.py}. Este script carrega o ficheiro de modelo guardado (\texttt{modelo\_ia.pkl}) e calcula o preço estimado para as características exatas introduzidas pelo utilizador (marca, modelo, quilometragem, potência, ano, combustível e tipo de transmissão), devolvendo o resultado para exibição imediata no ecrã.

A unificação dos dados é um pilar fundamental da arquitetura: todos os caminhos de leitura e escrita nas classes C++ e scripts de suporte foram alterados para apontar diretamente para os ficheiros centralizados localizados na pasta \texttt{JuOLXana} (\texttt{olx\_carros.csv}, \texttt{olx\_carros.json} e \texttt{olx\_carros\_ia.json}). Isto garante que a aplicação desktop e o portal web operam sob a mesma base de dados.

\section{Considerações Técnicas e Limitações}
A principal limitação do sistema reside na fragilidade inerente ao web scraping: alterações ao HTML do OLX (seletores CSS, estrutura do DOM, mecanismos de proteção anti-bot) podem inutilizar os parsers sem aviso prévio. Do ponto de vista ético e legal, o sistema opera em modo de leitura e não simula transações nem acede a áreas autenticadas. O ritmo de pedidos respeita os limites razoáveis de cortesia para com o servidor de origem.

\subsection{Melhorias Efetuadas e Resolução de Limitações}
O desenvolvimento final do sistema permitiu mitigar por completo as três principais limitações identificadas na fase exploratória da Inteligência Artificial:
\begin{enumerate}
    \item \textbf{Ausência da variável potência:} O scraper de recolha de dados foi atualizado para extrair dinamicamente a potência (\texttt{potencia\_cv}) diretamente da página de detalhes de cada anúncio. Adicionalmente, implementou-se uma heurística de estimação de potência com base em termos frequentes do título (ex. ``320d'', ``1.6 tdi'') e médias de marca/modelo para preencher valores em falta. A inclusão da potência como feature activa resultou num aumento drástico do $R^2$ de $\sim 0.35$ para mais de $0.60$ nos principais modelos regressivos.
    \item \textbf{Falta de persistência do modelo:} Criou-se um mecanismo de gravação que exporta o estado completo do melhor modelo selecionado para o ficheiro \texttt{modelo\_ia.pkl}. O estado persistido inclui os pesos do regressor, o mapeamento de variáveis categóricas frequentes e a normalização de features numéricas (\texttt{StandardScaler}).
    \item \textbf{Problemas de importação de classes na deserialização:} Para evitar falhas ao importar classes customizadas entre diferentes pastas (\texttt{GeOLXavani} e \texttt{JuOLXana}), o ficheiro pickle passou a armazenar apenas estruturas nativas de Python e do scikit-learn/xgboost. O script de avaliação (\texttt{evaluate\_car.py}) reconstrói a matriz de features e as dummy variables dinamicamente a partir deste dicionário de estado, garantindo total desacoplamento e estabilidade.
\end{enumerate}

\section{Conclusão}
Este trabalho demonstrou o desenvolvimento de um produto de dados completo, do princípio ao fim. Foi comprovado que, com ferramentas de código aberto (Python, scikit-learn, n8n) e de desenvolvimento nativo (Qt C++), é possível transcender a simples extração de informação, incorporando inteligência analítica baseada em Machine Learning para resolver assimetrias de informação no mercado automóvel. A orquestração do scraper com os algoritmos preditivos fornece métricas acionáveis, prontamente comunicadas ao utilizador através do portal web, aplicação desktop e das notificações automáticas.

Como trabalho futuro, identifica-se a integração de modelos de Natural Language Processing (NLP) para análise semântica das descrições dos vendedores, permitindo a deteção automática de sinistros ou extras de luxo; a integração com outras fontes de dados (Standvirtual, CustoJusto); e a migração para uma stack com base de dados persistente (PostgreSQL + API REST) que permita armazenamento histórico, multi-utilizador e otimização contínua dos modelos preditivos.

\begin{thebibliography}{9}
\bibitem{beauty} Richardson L (2007) Beautiful Soup Documentation. Crummy.com. \url{https://www.crummy.com/software/BeautifulSoup/bs4/doc/}
\bibitem{selenium} Selenium Project (2024) Selenium WebDriver Documentation. \url{https://www.selenium.dev/documentation/}
\bibitem{n8n} n8n GmbH (2024) n8n Documentation --- Workflow Automation Tool. \url{https://docs.n8n.io/}
\bibitem{olx} OLX Portugal (2025) Carros, Motos e Barcos --- Anúncios. \url{https://www.olx.pt/carros-motos-e-barcos/carros/}
\bibitem{pandas} McKinney W (2017) Python for Data Analysis, 2nd edn. O'Reilly Media, Sebastopol.
\bibitem{xgboost} Chen T, Guestrin C (2016) XGBoost: A Scalable Tree Boosting System. KDD '16. \url{https://doi.org/10.1145/2939672.2939785}
\bibitem{lgbm} Ke G et al. (2017) LightGBM: A Highly Efficient Gradient Boosting Decision Tree. NeurIPS 2017.
\end{thebibliography}

\end{document}
